\section{Evaluation and Conclusion}
\label{sec:evaluation-conclusion}

\graphicspath{{../evidence/}{evidence/}}

Chapter~9 closed the experimental sequence by showing that the two tested Copy
Fail implementations no longer produced the unauthorised privilege transition
on the corrected Ubuntu kernel. This final chapter does not introduce another
exploit run or another defensive prototype. Its purpose is to evaluate the
complete chain of evidence, draw the six research-question answers together,
and state exactly what the project demonstrated.

The investigation moved through five materially different security states. It
first established a vulnerable runtime and reproduced the UID~0 transition.
It then observed and detected the resulting privileged-file view, evaluated two
generations of a custom compensating control, compared those controls with
Canonical's temporary module-loading mitigation, and finally retested after the
vendor kernel correction. Treating these phases as one sequence matters because
the same visible outcome can have different meanings. A failed proof of concept
may indicate a policy denial, a missing interface, a broken sample, or a repaired
kernel. The controls surrounding each phase determine which conclusion the
evidence can support.

The overall result is positive, but intentionally asymmetric. Reproduction was
demonstrated for a specific reconstructed Ubuntu environment. Detection used a
behavioural integrity signal whose decision did not depend on the proof-of-
concept identity, but that signal did not uniquely attribute Copy Fail. MORI
v2.7 selectively denied the two tested same-process exploitation paths while
preserving the exercised interface operations, but it retained explicit bypass,
lifecycle, and fail-open boundaries. The corrected vendor kernel produced the
strongest operational result because it removed the reproduced behaviour while
leaving the tested AF\_ALG interface operations available.

\begin{evidencebox}
The final synthesis is anchored to the evidence already evaluated in the
experimental chapters. The decisive checkpoints are the vulnerable UID~0
transition and read-path divergence
(\cref{fig:phase02-uid0-transition,fig:phase02-view-divergence}), MORI's
behavioural alert (\cref{fig:phase03-mori-detection-root}), selective denial of
the Python and independent C paths
(\cref{fig:phase04-v27-python-denied,fig:phase04-v27-c-denied}), the vendor and
custom module-policy captures
(\cref{fig:phase01-vendor-kmod-reference,fig:phase04-v1-guard-config-comparison}),
and the patched-runtime regression results
(\cref{fig:phase05-python-poc-regression,fig:phase05-c-poc-regression,fig:phase05-final-target-integrity}).
No new experimental evidence is introduced in this chapter.
\end{evidencebox}

\subsection{Consolidated Findings}

The six research questions form a dependency chain rather than six unrelated
pass or fail checks. RQ1 establishes that there is a real laboratory effect to
study. RQ2 and RQ3 ask what can be observed without depending on one published
sample. RQ4 and RQ5 move from observation to temporary interruption and compare
the cost of selective and broad controls. RQ6 then tests the intended final
state: a corrected kernel without continued denial of the affected interface.

\Cref{tab:final-rq-synthesis} summarises the supported answer to each question
and the boundary that must remain attached to it.

\begin{table}[H]
    \centering
    \caption{Consolidated answers to the six research questions.}
    \label{tab:final-rq-synthesis}
    \scriptsize
    \begin{tabularx}{\textwidth}{@{}>{\bfseries\raggedright\arraybackslash}p{0.09\textwidth}>{\raggedright\arraybackslash}p{0.44\textwidth}Y@{}}
        \toprule
        Question & Supported answer & Principal boundary \\
        \midrule
        RQ1 &
        Yes. On Ubuntu kernel \code{6.8.0-116-generic}, a command launched from
        the recorded UID~1001 \code{labuser} context produced a UID~0 shell. &
        Operational repeatability was observed on one reconstructed VM. The
        preserved vulnerable launch has no immediate pre-launch sample hash,
        and no cross-host or statistical trial matrix was collected. \\

        RQ2 &
        The changed mounted view of \code{/usr/bin/su} was the most consistent
        directly observed detection opportunity; kernel and module state added
        exposure context, while UID~0 confirmed impact. &
        No complete syscall, eBPF, page-cache, or writeback trace validated the
        entire lower-level exploit chain. \\

        RQ3 &
        Yes. MORI detected deviation from a known-good privileged target using
        logic that did not match the proof-of-concept filename, language, or
        exact sample hash. &
        The digest mismatch was cause-agnostic, target-bounded, and observed
        after the changed view existed. It did not uniquely identify Copy Fail
        or measure time to alert. \\

        RQ4 &
        Yes, within the tested model. MORI v2.7 denied the preserved Python and
        independent C same-process paths before privilege escalation while the
        exercised benign interface operations remained available. &
        The result does not cover every process split, lifecycle edge case,
        allocation failure, target, workload, or future exploit variation. \\

        RQ5 &
        Canonical's \code{kmod} rule is the stronger broad emergency measure
        when interface removal is acceptable. MORI v2.7 demonstrated narrower
        experimental prevention where tested interface availability mattered. &
        MORI adds privileged local code, maintenance, kernel dependencies, and
        assurance obligations that the packaged vendor mitigation avoids. \\

        RQ6 &
        Yes, within the captured retest. On kernel
        \code{6.8.0-137-generic}, the Python path no longer produced its prior
        UID~0 outcome, the independent C path also remained unprivileged, and
        the exercised AF\_ALG operations remained reachable. &
        The identities differed between the vulnerable and patched Python runs,
        the C path lacks a preserved successful unmitigated vulnerable run, and
        the benign control did not perform encryption or decryption. \\
        \bottomrule
    \end{tabularx}
\end{table}

Taken together, these answers establish more than a successful exploit and a
subsequent failed retest. The vulnerable-side evidence records an ordinary-user
starting state, the direct UID~0 result, the altered mounted view, the divergent
backing-filesystem read, and recovery after reboot
(\cref{fig:phase02-pre-exploit,fig:phase02-uid0-transition,fig:phase02-view-divergence,fig:phase02-reboot-recovery}).
The defensive phases then show three distinct responses to that path: observing
its consequence, denying selected prerequisites or correlations, and removing
the underlying vulnerable behaviour through the kernel correction.

The differences between those responses are security-relevant. Detection can
support investigation without preventing impact. A broad temporary mitigation
can prevent the known path by removing functionality. A selective compensating
control can preserve more functionality, but only by accepting more policy
complexity and a narrower assurance claim. A fixed kernel addresses the defect
at its source and therefore remains the appropriate end state.

\subsection{Detection Limitations}

The detection work produced a useful behavioural result, but not a universal
Copy Fail signature. MORI's integrity monitor compared protected privileged
targets with a known-good baseline and detected the changed mounted view of
\code{/usr/bin/su}. Its decision logic did not inspect the proof-of-concept
filename, language, or sample hash. The exact-hash YARA comparison, by contrast,
stopped matching after a harmless one-byte variation
(\cref{fig:phase03-yara-known-match,fig:phase03-yara-variant-no-match}). This
supports the architectural distinction between behavioural and exact-sample
detection. It does not demonstrate equal detection coverage for every possible
implementation.

It did not generalise to unique vulnerability attribution. A changed digest can
also be caused by an authorised package update, corruption, an administrative
action, or a different attack technique. The detector answered that a protected
object no longer matched its baseline; it did not prove why the mismatch
occurred. The vulnerable kernel version, loaded-module state, execution context,
and UID transition can strengthen an investigation, but the project did not
validate one complete real-time correlation joining all of those signals.

The timing boundary is equally important. The monitored integrity deviation was
visible after the changed target view existed. The later UID~0 result confirmed
impact, but an alert produced after privilege escalation cannot be presented as
prevention. Polling also creates a sampling boundary: a deviation that appears
and disappears between complete scans may not be observed. The configured
sleep interval was an implementation setting, not a measured end-to-end alert
latency, and the experiment therefore reports no time-to-alert metric.

The detection evaluation was further limited by coverage and trust:

\begin{itemize}
    \item only \code{/usr/bin/su} was experimentally demonstrated as the Copy
    Fail target; the remaining baseline entries were coverage candidates;

    \item the baseline did not include every privileged executable, file
    capability, arbitrary target path, or container host boundary;

    \item no representative benign-workload corpus, false-positive rate,
    sustained performance benchmark, or multi-host deployment was measured;

    \item the monitor, baseline, and local journal could all become untrustworthy
    after unrestricted root compromise unless measurements are protected or
    exported independently; and

    \item the proposed lower-level AF\_ALG and \code{splice} signals became
    inputs to the later prevention design, but were not independently preserved
    as a complete detection trace during the original monitored exploit run.
\end{itemize}

The practical role of the detector is therefore bounded but still valuable. It
can reveal a security-relevant privileged-target deviation and supply evidence
for response, especially when combined with exposure and identity context. It
should not be marketed as proof of Copy Fail, as a guarantee of timely alerting,
or as a substitute for prevention and remediation.

\subsection{Mitigation Limitations}

The mitigation phases show that successful blocking is not a complete security
evaluation. MORI Guard v1 and Canonical's temporary mitigation both interrupted
the known module-dependent path by making \code{algif\_aead} unavailable through
the tested loading route. That simplicity provides a broad containment boundary,
but it also removes benign access to the affected interface. MORI Guard v1 added
local telemetry without improving that compatibility trade-off, which left the
vendor-packaged rule as the stronger operational choice for emergency use.

MORI v2.7 pursued a different goal. It retained the exercised AF\_ALG socket,
bind, and accept operations and denied the correlated SUID-to-pipe operation for
the Python and independent C paths. The same suite also showed that individual
signals, an alternate target, a cross-TGID sequence, and the monitor's own
expected reads were not indiscriminately denied. This demonstrated selectivity
relative to v1, not universal compatibility
(\cref{fig:phase04-v27-python-denied,fig:phase04-v27-python-post-state,fig:phase04-v27-c-denied,fig:phase04-v22-benign-bind}).

The price of that selectivity is a larger assumption surface. The policy depends
on BPF LSM and BTF/CO-RE support, privileged deployment, correct process and
thread lifecycle tracking, successful state-map updates, and continued trust in
the loader and monitor. Process separation was intentionally allowed and remains
an evasion boundary. Matching benign software may be denied, state-allocation
failure is fail-open, and fork, exec, descriptor transfer, PID reuse,
multithreaded races, and abnormal teardown were not exhaustively evaluated.
The targeted regression suite improves confidence in the frozen candidate but
does not remove those architectural limits.

The four defensive states can therefore be ordered by purpose rather than by a
single claim that one control is always better:

\begin{table}[H]
    \centering
    \caption{Operational interpretation of the evaluated defensive states.}
    \label{tab:final-control-interpretation}
    \small
    \begin{tabularx}{\textwidth}{@{}>{\bfseries\raggedright\arraybackslash}p{0.22\textwidth}>{\raggedright\arraybackslash}p{0.31\textwidth}Y@{}}
        \toprule
        State & Appropriate interpretation & Required caution \\
        \midrule
        MORI detection &
        Host-local observation and evidence of privileged-target deviation. &
        Alerts after the effect and does not uniquely attribute the cause. \\

        MORI Guard v1 / Canonical rule &
        Broad temporary containment by removing the module prerequisite. &
        Benign interface access is also removed; the vendor rule is easier to
        support operationally. \\

        MORI v2.7 &
        Selective experimental denial for the tested behavioural correlation. &
        Narrow assurance, custom privileged code, explicit evasion and
        fail-open boundaries, and no production workload validation. \\

        Corrected vendor kernel &
        Long-term remediation of the reproduced vulnerable behaviour while
        retaining the tested interface operations. &
        Deployment still requires reboot, version verification, and a bounded
        regression test; one laboratory cannot prove every variant or platform. \\
        \bottomrule
    \end{tabularx}
\end{table}

Canonical identifies \code{6.8.0-117.117} as the first corrected Noble generic
kernel, and the laboratory's final \code{6.8.0-137.137} runtime was beyond that
boundary \cite{ubuntu-copyfail-cve}. Its observed behaviour also matched the
expected consequence of the upstream return to out-of-place AEAD processing
\cite{linux-copyfail-fix}. That evidence supports the fixed kernel as the final
remediation without turning the laboratory result into a claim about every
distribution, kernel flavour, target, or exploit implementation.

\subsection{Lessons Learned and Future Work}

The most useful lessons came from the places where the initial interpretation
or control was not good enough. A changed privileged-file digest was important,
but the direct UID~0 identity check was what proved successful exploitation.
An exact-sample rule was easy to build, but a one-byte variation showed how
little assurance it provided. MORI Guard v1 blocked the proof of concept, but
the benign control exposed its compatibility cost. The first selective designs
then produced their own false positives and lifecycle bypasses, forcing the
policy to be evaluated as software rather than celebrated as soon as one attack
failed.

\subsubsection{Evidence quality determines claim strength}

The evidence workflow was strongest where independent observations converged.
Examples include the ordinary-user baseline followed by direct UID~0 output,
the mounted-view and backing-filesystem digest comparison, the paired positive
and negative MORI controls, and the patched-side combination of interface
availability with absent privilege transition. Captured sample hashes add
provenance where they exist, while the missing immediate vulnerable-side Python
hash and the failed machine-readable C capture remain explicit gaps. Phase
separation, VM snapshots, manifests, and target-state checks made it possible to
relate those observations without treating every screenshot as self-sufficient
proof.

The workflow was weaker where provenance depended on filenames, documentation,
or inferred chronology. The misnamed Phase~05 baseline, failed machine-readable
C provenance capture, missing embedded timestamps, changed Python execution
identity, and screenshot-only retests do not invalidate the observations they
contain. They do limit causal precision. Preserving those imperfections in the
report is more useful than silently repairing the narrative, because it keeps
the conclusion aligned with what the artifacts can actually establish.

This leads to a general methodological lesson: evidence collection is part of
the experiment, not administrative work performed after it. Commands that show
identity, version, control isolation, sample provenance, execution, and final
target state should be captured in one planned sequence wherever possible.
When that sequence breaks, the appropriate response is to narrow the claim or
collect a clearly labelled supplement, not to reconstruct missing proof from
memory.

\subsubsection{Regression testing is part of defensive policy design}

MORI's development history shows why one successful denial is not enough. The
broad first control prevented exploitation and also removed benign
functionality. A later monitor-read false positive showed that a detector can
interfere with its own observation path. The retained 12-second regression then
demonstrated that policy state which looks correct during a short test may fail
after a lifecycle transition. Each failure changed the design and became part
of the final assurance boundary.

The positive exploit tests were therefore only half of the evaluation. The
negative controls, alternate target, cross-process case, trusted-monitor check,
state-expiry regression, and post-test target validation determined whether the
control was selective and whether denial actually preserved the security
objective. Future work on any similar behavioural control should begin with
that regression mindset rather than adding it after the first bypass appears.

\subsubsection{Priorities for future work}

The next useful experiments are those that close a specific evidence or
assurance gap:

\begin{enumerate}
    \item \textbf{Produce a matched evidence supplement.} Repeat vulnerable and
    patched tests with the same dedicated unprivileged account, timestamped
    isolation immediately before execution, machine-readable sample provenance,
    raw transcripts, boot identifiers, and matched mounted and backing-device
    checks. Retain the original imperfect artifacts rather than replacing them.

    \item \textbf{Expand the benign AF\_ALG control.} Perform real encrypt,
    decrypt, authentication-tag verification, concurrency, and application-level
    tests. The current control establishes tested interface availability, not a
    complete cryptographic compatibility result.

    \item \textbf{Stress the MORI state model.} Evaluate fork and exec paths,
    descriptor passing, process separation, multithreaded races, abnormal
    teardown, PID and TGID reuse, map pressure, allocation failure, and monitor
    lifecycle changes. Compare fail-open and fail-closed choices explicitly.

    \item \textbf{Measure operational cost.} Use representative benign
    workloads to quantify false positives, event loss, CPU and memory overhead,
    and decision latency across supported kernels, distributions, and
    architectures. Export telemetry to a separately trusted system before
    evaluating incident-response value.

    \item \textbf{Broaden without overgeneralising.} Test additional privileged
    targets, file-capability executables, kernel flavours, and hardware. A
    container-escape claim would require a separate threat model and dedicated
    reproduction rather than being inferred from this local privilege-escalation
    study.

    \item \textbf{Retire temporary controls after remediation.} After
    validating the fixed kernel, remove obsolete module or BPF policies and
    recheck the applications that depended on the affected interface.
\end{enumerate}

These priorities identify how future work could broaden the present laboratory
evidence. Until then, MORI v2.7 remains a research-grade example of selective
behavioural prevention, not a production recommendation.

\subsection{Conclusion}

This project asked more than whether a published local privilege-escalation
technique could print a root prompt. The laboratory reconstructed the vulnerable
Ubuntu runtime, recorded a UID~1001 starting state, and directly validated the
transition to UID~0. It then separated that result from the mounted-view anomaly:
the normal VFS path returned changed content, a direct ext4 read retained the
original digest, and reboot restored the mounted view.

The observation phase turned that effect into a behavioural integrity signal.
MORI detected the target deviation without binding its decision to one filename
or hash, but it did not provide unique attribution or prevention. Broad module
denial then interrupted the known path at the cost of the tested interface.
MORI v2.7 preserved the exercised operations and denied two correlated paths,
but only within its documented process, lifecycle, and deployment model.

The vendor comparison kept the custom result in its proper role. Canonical's
packaged module rule remained preferable for broad emergency containment when
temporary interface loss was acceptable, and neither temporary control replaced
the corrected kernel. On the patched runtime, the affected module loaded and
the exercised socket, bind, and accept operations remained reachable
(\cref{fig:phase05-patched-runtime-isolation,fig:phase05-benign-aead-control}).
The Python path did not reproduce its documented vulnerable-side UID~0 outcome,
the independent C path also remained unprivileged, and \code{/usr/bin/su}
retained its known-good captured state
(\cref{fig:phase05-python-poc-regression,fig:phase05-c-poc-regression,fig:phase05-final-target-integrity}).
That last observation establishes the target state at capture time, not a
continuous integrity guarantee.

The final conclusion is therefore:

\begin{lstlisting}
Copy Fail was reproducible as a local privilege escalation on the
tested vulnerable Ubuntu kernel and observable through a repeatable
privileged-target integrity deviation.

Behavioural detection and selective temporary prevention were possible
for the demonstrated path, but neither supplied unique attribution nor
universal protection.

The corrected vendor kernel removed the reproduced Python-path behaviour,
while the independent C path also remained unprivileged. The tested
AF_ALG socket, bind, and accept operations remained available, making
the vendor kernel the appropriate final remediation for this system.
\end{lstlisting}

That statement is deliberately narrow. The project used one principal vulnerable
VM, one Ubuntu kernel line, one documented Python reproduction, an independent C
implementation in later tests, a bounded target set, and a finite regression
suite. The C path was not preserved succeeding unmitigated on the vulnerable
kernel. No container escape, universal exploit coverage, population-level
success rate, or production readiness for MORI was demonstrated.

The lasting output is an evidence-linked path from demonstrated impact through
observation and temporary control to vendor remediation, with the remaining
uncertainty stated openly.